\documentclass[11pt,letterpaper]{article}
\usepackage[margin=1in,headheight=14pt]{geometry}
\usepackage[T1]{fontenc}
\usepackage{lmodern,microtype}
\usepackage{amsmath,amssymb,amsthm,mathtools,bm}
\usepackage{booktabs,longtable,array}
\usepackage{xcolor}
\definecolor{bcblue}{RGB}{25,65,105}
\usepackage[colorlinks=true,linkcolor=bcblue,citecolor=bcblue,urlcolor=bcblue,breaklinks=true]{hyperref}
\usepackage{bookmark,fancyhdr,needspace,etoolbox}
\pagestyle{fancy}
\fancyhf{}\fancyhead[L]{\small Beta mixtures across integer twists}
\fancyhead[R]{\small ProveIt research continuation}
\fancyfoot[C]{\thepage}
\setlength{\emergencystretch}{2em}
\setlength{\parskip}{3pt}
\allowdisplaybreaks[2]
\numberwithin{equation}{section}
\newtheorem{theorem}{Theorem}[section]
\newtheorem{proposition}[theorem]{Proposition}
\newtheorem{corollary}[theorem]{Corollary}
\newtheorem{lemma}[theorem]{Lemma}
\theoremstyle{definition}
\newtheorem{definition}[theorem]{Definition}
\newtheorem{example}[theorem]{Example}
\theoremstyle{remark}
\newtheorem{remark}[theorem]{Remark}
\BeforeBeginEnvironment{theorem}{\Needspace{6\baselineskip}}
\BeforeBeginEnvironment{proposition}{\Needspace{6\baselineskip}}
\BeforeBeginEnvironment{corollary}{\Needspace{6\baselineskip}}
\BeforeBeginEnvironment{lemma}{\Needspace{6\baselineskip}}
\newcommand{\bcFP}{\operatorname{FP}}
\newcommand{\bcPf}{\operatorname{Pf}}
\newcommand{\bcPV}{\operatorname{PV}}
\newcommand{\bcLi}{\operatorname{Li}}
\newcommand{\bcRes}{\operatorname{Res}}
\newcommand{\bcCT}{\operatorname{CT}}
\newcommand{\bcRe}{\operatorname{Re}}
\newcommand{\bcsgn}{\operatorname{sgn}}
\newcommand{\bcC}{\mathcal C}
\newcommand{\bcY}{\mathcal Y}
\newcommand{\bcZ}{\mathcal Z}
\newcommand{\bcJ}{\mathcal J}
\newcommand{\bcT}{\mathbb T}
\newcommand{\bcN}{\mathbb N}
\newcommand{\bcR}{\mathbb R}
\newcommand{\bcD}{\mathcal D}
\newcommand{\bce}{\mathrm e}
\DeclareMathOperator{\rank}{rank}
\hypersetup{pdftitle={Beta Mixtures Across Integer Twists},pdfauthor={Research continuation prepared for ProveIt},pdfsubject={Finite spectral corrections, harmonic moments, polylogarithms, Hurwitz zeta and Stieltjes primitives}}
\title{\textbf{Beta Mixtures Across Integer Twists}\\[5pt]
\large Finite spectral corrections, harmonic logarithmic moments,\\
and Stieltjes--polylogarithm primitives}
\author{Research continuation prepared for the ProveIt project}
\date{October 10, 2026\\[3pt]\small Reference commit: \texttt{446dca9bd65c20e409d1d2561d4ffad408103e2a}}
\begin{document}
\maketitle
\begin{abstract}
The ProveIt manuscript evaluates convolutions of unequal twisted Hurwitz
kernels by Dirichlet beta mixtures when the twists lie in one interval
between consecutive integers. We determine the correction when that
hypothesis is removed and the total order is a positive integer. For any
finite number of factors, the difference between the convolution and the
affine finite-part mixture is an explicit finite Fourier polynomial. Its
coefficient at a crossed mode is a sine of the sum of the orders on the
negative side. Although this correction vanishes at every positive-integer
order tuple, its balanced order derivatives need not vanish. We give all
such derivatives and prove that their first-order rank is one less than
the number of occupied integer chambers.

A two-endpoint specialization evaluates every sufficiently vanishing
centered-logarithm moment by a polynomial times $\pi u\cot\pi u$. At integer
endpoint orders the resulting ordinary integrals are finite rational
linear combinations of even zeta values. Their parameter primitives have
finite classical-polylogarithm expressions, producing convergent
inverse-logarithm integrals with odd zeta values. An independent logistic
Mellin transform gives finite Hurwitz-zeta reductions, all logarithmic
finite parts, and generalized Stieltjes antiderivatives. In particular,
the finite part of an unnormalized Mellin integral requires an explicit
Gamma-derivative contribution from the pole residue.

The proofs are self-contained apart from standard special-function
identities, which are also derived where used. The package records 508
finite exact checks, including three negative controls, and 94
floating-point diagnostics. These checks are not a formalization or
interval certification. The contribution is an extension of the inspected
repository material, not a claim of literature priority.
\end{abstract}
\newpage
\begingroup
\setlength{\parskip}{0pt}
\tableofcontents
\endgroup
\newpage

\section{Research target, relationship to the manuscript, and scope}
\label{bc:sec:introduction}

The immediate starting point is the manuscript section
\emph{Unequal twists through beta mixtures and finite resolvent identities},
file \texttt{07-mixed-twist-identities.tex} \cite{repo-mixed}. Its kernel
$K_\alpha^\theta$ has Fourier multiplier $(2\pi i(n+\theta))^{-\alpha}$.
When all real twists belong to a single open interval between consecutive
integers, its Dirichlet-mixture theorem is valid for arbitrary orders with
positive real parts. The same section gives finite partial-fraction
reductions at positive integer orders without the common-interval
restriction. Both statements are correct as written.

There is a genuine gap between those regimes: the beta simplex can cross
an integral twist even when none of its vertices is integral. The
intermediate kernel then has a pole in the twist parameter. Simply
replacing the integral by a principal value, or by the real finite part at
higher order, does not in general preserve the identity. Moreover, an
identity valid only on the discrete set of integer order tuples cannot be
differentiated in the orders without first constructing its analytic
extension.

The main theorem below fills this gap at every positive integer
\emph{total} order, allowing the individual orders to be fractional. It is
not restricted to two factors. The resulting correction is especially
simple because only finitely many Fourier modes can change sign. This
finite-dimensional feature makes all balanced order derivatives explicit
and permits a complete rank calculation. The same scalar computation
produces a family of exact ordinary integrals, rather than estimates.

\subsection{What is established here}

The central results are Theorem~\ref{bc:thm:periodic} (arbitrary-depth
integer-total crossing formula), Theorem~\ref{bc:thm:jets} (all balanced
order derivatives), and Theorem~\ref{bc:thm:rank} (the complete
chamberwise criterion for their vanishing). The analytic consequences are
Theorem~\ref{bc:thm:moments}, the polylogarithm primitive in
Proposition~\ref{bc:prop:polyprimitive}, the logarithmic Mellin finite-part
formula in Theorem~\ref{bc:thm:mellin-fp}, and the Stieltjes primitive ladder
in Theorem~\ref{bc:thm:primitives}. Every one is proved below.

The beta integral, Gamma reflection, integer polygamma--Hurwitz
conversion, elementary polylogarithm differential rule, and finite
Hurwitz expansion of the equal-period Barnes-type zeta are classical
ingredients \cite{dlmf-beta,dlmf-gamma,dlmf-polygamma,dlmf-polylog,dlmf-hurwitz}.
Their use is not itself a novelty claim. The distinction between analytic
boundary values and uncorrected Feynman-parameter integrals also has a
substantial history; the papers of Weldon and Bedaque--Das provide relevant
context \cite{weldon,bedaque-das}. Only their abstracts were consulted for
this contextual comparison. No step of our proof is outsourced to those
papers, and no assertion is made that the general boundary-value principle
originates here.

\subsection{What is not claimed}

This is not a proof of the manuscript's $S_6$ or revised $S_8$ period
conjectures. The repository README continues to mark them conjectural
\cite{repo-readme}. Nor do we claim a finite reduction of all arbitrary
unequal-twist order derivatives to a prescribed list of single Hurwitz
values. Our rank theorem concerns explicit Fourier corrections, not
arithmetic independence of periods.

The targeted source was read again at the reference commit printed on
the title page. The manuscript README, the chapter inventory, and the
incoming directory listing were inspected. The harmonic Gamma-coefficient
section was read in full; it already proves the corrected all-order Choi
families \cite{repo-harmonic}. Only the beginning of the log-Gamma-moments
section was inspected. The six binary archives present in the incoming
listing were \emph{not} opened. Consequently this is neither an exhaustive
repository audit nor an exhaustive comparison with all pending reports.
The machine-readable provenance record makes these boundaries explicit.

% INTEGRATION-BEGIN
\section{Fourier, branch, and finite-part conventions}
\label{bc:sec:conventions}

Write $\bcT=\bcR/\mathbb Z$. For a periodic distribution $f$, use
\[
 \widehat f(n)=\langle f,\bce^{-2\pi i n x}\rangle,
 \qquad f=\sum_{n\in\mathbb Z}\widehat f(n)\bce^{2\pi i n x}.
\]
The pairing agrees with integration over $[0,1]$ for integrable functions.
Convolution is normalized so that Fourier coefficients multiply. For a
real nonintegral $\theta$, define the entire distribution-valued family
\begin{equation}\label{bc:eq:kernel}
 \widehat K_\alpha^\theta(n)
  =\exp\left[-\alpha\left\{
       \log(2\pi|n+\theta|)
       +\frac{\pi i}{2}\bcsgn(n+\theta)\right\}\right].
\end{equation}
On each compact set of $\alpha$, these coefficients and their order
derivatives have a common polynomial growth bound in $n$. Pairing against
a smooth test function therefore proves normal convergence and entire
dependence on $\alpha$. The same argument defines convolution of any
finite collection of these kernels.

For positive integers $m$ and $0<x<1$,
\begin{equation}\label{bc:eq:integer-kernel}
 K_1^\theta(x)=\frac{\bce^{-2\pi i\theta x}}{1-\bce^{-2\pi i\theta}},
 \qquad
 K_m^\theta(x)=
 \frac{(-1)^{m-1}}{(2\pi i)^{m-1}(m-1)!}
          \partial_\theta^{m-1}K_1^\theta(x).
\end{equation}
Indeed the first expression has Fourier coefficient
$(2\pi i(n+\theta))^{-1}$, by a direct integral. Differentiating this
coefficient gives the second expression. Endpoint representatives are
not assigned by this formula; the periodic distribution is the primary
object.

\subsection{The affine finite part}

For a positive integer $m$ we use the distributional identity
\begin{equation}\label{bc:eq:boundary}
 (y-i0)^{-m}
 =\bcPf y^{-m}
  +\frac{i\pi(-1)^{m-1}}{(m-1)!}\delta^{(m-1)}(y).
\end{equation}
Here $\bcPf y^{-1}$ is the Cauchy principal value, and
\[
 \bcPf y^{-m}=\frac{(-1)^{m-1}}{(m-1)!}
                     \frac{d^{m-1}}{dy^{m-1}}\bcPV\frac1y.
\]
For completeness, write $(y-i\varepsilon)^{-1}$ as
$y/(y^2+\varepsilon^2)+i\varepsilon/(y^2+\varepsilon^2)$.
The first term tends to the principal value and the second to $i\pi\delta$.
Differentiating $m-1$ times proves \eqref{bc:eq:boundary}. Replacing $i0$
by $-i0$ reverses the sign of the delta term.

If $h$ is smooth near zero and integrable away from zero on a compact
interval, write
\[
 \bcFP\int h(y)y^{-m}\,dy=\langle\bcPf y^{-m},h\rangle.
\]
We also allow a finite measure away from zero; only its restriction near
zero needs a smooth density. A partition of unity reduces this definition
to a smooth function of compact support near zero and an ordinary integral
against the measure elsewhere. An explicit
calculation uses Taylor subtraction through degree $m-1$ near zero,
followed by the exact finite parts of the subtracted monomials. Thus the
notation specifies a computable normalization, not merely a rule to
``discard infinity.'' Pairing \eqref{bc:eq:boundary} with $h$ gives
\begin{equation}\label{bc:eq:boundary-pair}
 \int h(y)(y-i0)^{-m}\,dy
 =\bcFP\int h(y)y^{-m}\,dy
       +\frac{i\pi}{(m-1)!}h^{(m-1)}(0).
\end{equation}

All simplex finite parts below use the displayed \emph{affine scalar}
$y=\sum t_j y_j$. This choice matters: an arbitrary nonlinear change of
variable can alter local finite-part terms. Later, logarithmic Mellin
finite parts use a different, explicitly stated cutoff variable. We
change variables between the two settings only for ordinary convergent
integrals, unless a separate normalization is supplied.

\section{A scalar Dirichlet crossing identity at every depth}
\label{bc:sec:scalar}

Let
\[
 \Delta_{d-1}=\{(t_1,\ldots,t_d):t_j\ge0,\ \sum_jt_j=1\},
 \qquad dt=dt_1\cdots dt_{d-1},
\]
where $t_d=1-\sum_{j<d}t_j$. Fix nonzero real numbers $y_j$ and complex
parameters $\alpha_j$ satisfying
\begin{equation}\label{bc:eq:scalar-assumptions}
 \bcRe\alpha_j>0,\qquad \sum_{j=1}^d\alpha_j=m\in\bcN.
\end{equation}
Set
\begin{equation}\label{bc:eq:scalar-data}
 \sigma=\sum_{y_j<0}\alpha_j,
 \qquad Q=\prod_{j=1}^d|y_j|^{-\alpha_j}.
\end{equation}
The normalized pushforward $h$ is specified by
\begin{equation}\label{bc:eq:pushforward}
 \int h(y)F(y)\,dy=
 \frac{\Gamma(m)}{\prod_j\Gamma(\alpha_j)}
 \int_{\Delta_{d-1}}\prod_jt_j^{\alpha_j-1}
                          F\left(\sum_jt_jy_j\right)dt.
\end{equation}
For complex parameters this is a complex measure, not a probability
density. Positivity is unnecessary.

\begin{lemma}[Smoothness at an interior crossing]\label{bc:lem:smooth}
If the $y_j$ include both signs, the pushforward in
\eqref{bc:eq:pushforward} has a smooth density on a neighborhood of zero.
Every fixed derivative there depends holomorphically on the parameters
in \eqref{bc:eq:scalar-assumptions}. If all signs agree, its support is
bounded away from zero.
\end{lemma}
\begin{proof}
Group the negative and positive indices. The usual change of simplex
coordinates first chooses the total negative mass $v$, then the internal
coordinates in each group. It factors the normalized measure into two
internal Dirichlet measures and the beta density
\[
 \frac{\Gamma(m)}{\Gamma(\sigma)\Gamma(m-\sigma)}
               v^{\sigma-1}(1-v)^{m-\sigma-1}\,dv.
\]
This factorization follows directly by its Jacobian and the beta
integral; it remains valid for positive real parts by absolute
convergence. Conditional on the internal coordinates, let $Y_-$ and
$Y_+$ be the corresponding weighted means of the negative and positive
$y_j$. Some $c>0$ satisfies $Y_-\le-c$ and $Y_+\ge c$ uniformly.
Now $y=vY_-+(1-v)Y_+$, so
\[
 v=\frac{Y_+-y}{Y_+-Y_-},\qquad
 \left|\frac{dv}{dy}\right|=\frac1{Y_+-Y_-}.
\]
For $|y|<c/2$, both $v$ and $1-v$ stay uniformly away from zero, because
$Y_\pm$ also range in fixed compact sets. Every $y$ derivative of the
conditional beta density is therefore bounded by a constant times an
integrable internal Dirichlet weight. Differentiation under the integral
proves smoothness. The same domination on compact parameter sets, with
additional endpoint logarithms when needed, proves holomorphic
dependence. The one-sign assertion is immediate from convexity.
\end{proof}

\begin{theorem}[Finite part and crossing derivative]\label{bc:thm:scalar}
Under \eqref{bc:eq:scalar-assumptions},
\begin{align}
 \bcFP\int h(y)y^{-m}\,dy&=\cos(\pi\sigma)Q,
       \label{bc:eq:scalar-fp}\\
 h^{(m-1)}(0)&=\frac{(m-1)!}{\pi}\sin(\pi\sigma)Q.
       \label{bc:eq:density-derivative}
\end{align}
In the one-sign case the second formula is read as zero, since
$\sigma=0$ or $\sigma=m$. The two boundary values are
\begin{equation}\label{bc:eq:scalar-boundaries}
 \int h(y)(y\mp i0)^{-m}\,dy=\bce^{\pm i\pi\sigma}Q.
\end{equation}
\end{theorem}
\begin{proof}
For complex numbers $z_j$ in the right half-plane, the product of Gamma
integrals and the substitution $u_j=rt_j$ give
\begin{equation}\label{bc:eq:dirichlet}
 \prod_j z_j^{-\alpha_j}
 =\frac{\Gamma(m)}{\prod_j\Gamma(\alpha_j)}
 \int_{\Delta_{d-1}}\frac{\prod_jt_j^{\alpha_j-1}}
                              {(\sum_jt_jz_j)^m}\,dt.
\end{equation}
Indeed the Jacobian is $r^{d-1}$ and integration in $r$ produces
$\Gamma(m)(\sum t_jz_j)^{-m}$. This proves the identity directly in a
domain where all relevant integrals converge absolutely.

Put $z_j=\varepsilon+i y_j$ in \eqref{bc:eq:dirichlet} and factor out
$i^{-m}$ from both sides. As $\varepsilon\downarrow0$, the left side
becomes $Q\bce^{i\pi\sigma}$. For a negative $y_j$ the factor
$(y_j-i0)^{-\alpha_j}$ contributes $\bce^{i\pi\alpha_j}$; this is the
entire phase correction. On the right, Lemma~\ref{bc:lem:smooth} and
\eqref{bc:eq:boundary-pair} justify the distributional limit at zero;
away from zero dominated convergence applies to the finite pushforward
measure. This proves the upper sign of
\eqref{bc:eq:scalar-boundaries}. The substitution
$z_j=\varepsilon-i y_j$ proves the lower sign by the same argument.
Average the two identities to obtain \eqref{bc:eq:scalar-fp}; subtract
them and use \eqref{bc:eq:boundary-pair} to obtain
\eqref{bc:eq:density-derivative}. No reality assumption on the parameters
is needed in this subtraction.
\end{proof}

\begin{remark}[A direct two-point version]
For $A,B>0$, $\bcRe\alpha,\bcRe\beta>0$ and $\alpha+\beta=m$, the theorem
states
\begin{equation}\label{bc:eq:two-beta}
 \frac{\Gamma(m)}{\Gamma(\alpha)\Gamma(\beta)}
 \bcFP\int_0^1
 \frac{t^{\alpha-1}(1-t)^{\beta-1}}
      {[B-(A+B)t]^m}\,dt
 =\frac{\cos(\pi\alpha)}{A^\alpha B^\beta}.
\end{equation}
The missing imaginary boundary term is
$i\sin(\pi\alpha)/(A^\alpha B^\beta)$. Here the affine singularity is
$t=B/(A+B)$, not necessarily $1/2$.
\end{remark}

\section{The finite spectral correction for unequal twists}
\label{bc:sec:periodic}

For nonintegral real twists $\theta_1,\ldots,\theta_d$, define
\begin{equation}\label{bc:eq:crossed-modes}
 \bcC=\{n\in\mathbb Z:\min_j(n+\theta_j)<0<\max_j(n+\theta_j)\}.
\end{equation}
This is a finite set. For every mode put
\[
 \sigma_n=\sum_{n+\theta_j<0}\alpha_j,
 \qquad Q_n=\prod_j|n+\theta_j|^{-\alpha_j}.
\]
Let $\mathcal M_m(\bm\alpha;\bm\theta)$ denote the simplex mixture whose
$n$th Fourier coefficient is the affine finite part of
\[
 \frac1{(2\pi i)^m}
 \frac{\Gamma(m)}{\prod_j\Gamma(\alpha_j)}
 \int_{\Delta_{d-1}}\prod_jt_j^{\alpha_j-1}
                     \left(n+\sum_jt_j\theta_j\right)^{-m}dt.
\]
Equivalently it is the affine finite-part interpretation of the integral
of $K_m^{\sum t_j\theta_j}$. Only finitely many modes need subtraction.

\begin{theorem}[Arbitrary-depth integer-total crossing formula]
\label{bc:thm:periodic}
If $\bcRe\alpha_j>0$ and $\sum_j\alpha_j=m\in\bcN$, then
\begin{align}
 K_{\alpha_1}^{\theta_1}*\cdots*K_{\alpha_d}^{\theta_d}
 ={}&\mathcal M_m(\bm\alpha;\bm\theta)
 +\frac{i}{(2\pi i)^m}\sum_{n\in\bcC}
      \bce^{2\pi i n x}\sin(\pi\sigma_n)Q_n.
 \label{bc:eq:periodic-main}
\end{align}
The mixture and its correction are holomorphic in the indicated
parameter hyperplane. If all twists lie in one open unit interval, the
correction is empty. If every $\alpha_j$ is a positive integer, the
correction vanishes even when $\bcC$ is nonempty.
\end{theorem}
\begin{proof}
At mode $n$, formula \eqref{bc:eq:kernel} gives the product multiplier
\begin{equation}\label{bc:eq:mode-phase}
 \prod_j(2\pi i(n+\theta_j))^{-\alpha_j}
       =(2\pi i)^{-m}\bce^{i\pi\sigma_n}Q_n.
\end{equation}
Theorem~\ref{bc:thm:scalar} evaluates the mixture coefficient as
$(2\pi i)^{-m}\cos(\pi\sigma_n)Q_n$. Their difference is the displayed
sine term. For $n\notin\bcC$, $\sigma_n$ is either $0$ or $m$, so the
difference is zero. This also proves the assertions about the common
interval and integer tuples.

For $|n|$ sufficiently large all $n+\theta_j$ have the same sign and are
comparable to $|n|$. Uniformly on compact parameter sets in the
hyperplane, $Q_n=O(|n|^{-m})$. The finitely many crossed coefficients
are holomorphic by Lemma~\ref{bc:lem:smooth}. Pairing the series with a
smooth test function now justifies the mixture, all parameter
derivatives, and the equality of distributions by Fourier uniqueness.
In fact at $m\ge2$ the undifferentiated coefficients are absolutely
summable. The distributional formulation also covers $m=1$ without any
endpoint convention.
\end{proof}

The correction in \eqref{bc:eq:periodic-main} is a finite trigonometric
polynomial in $x$. The local delta derivative used to obtain it was in
the \emph{intermediate twist variable}. It must not be replaced by a
Dirac distribution in $x$.

\begin{corollary}[Two twists crossing several integers]
\label{bc:cor:two-twists}
For nonintegral $a<b$ and $\alpha+\beta=m$ with positive real parts,
\begin{align}
 K_\alpha^a*K_\beta^b
 ={}&\frac{\Gamma(m)}{\Gamma(\alpha)\Gamma(\beta)}
 \bcFP\int_0^1(1-t)^{\alpha-1}t^{\beta-1}
                            K_m^{(1-t)a+tb}\,dt\notag\\
 &+\frac{i\sin(\pi\alpha)}{(2\pi i)^m}
     \sum_{\substack{k\in\mathbb Z\\a<k<b}}
       \frac{\bce^{-2\pi i kx}}{(k-a)^\alpha(b-k)^\beta}.
 \label{bc:eq:two-twist-main}
\end{align}
\end{corollary}
\begin{proof}
The crossed modes are exactly $n=-k$ with $a<k<b$. Their negative group
consists of the first factor, so $\sigma_{-k}=\alpha$.
\end{proof}

\subsection{Exact counterexamples to an uncorrected extension}

At $a=-1/2$, $b=1/2$ and $\alpha=\beta=1/2$, only $n=0$ crosses. Hence
\begin{equation}\label{bc:eq:half-counterexample}
 K_{1/2}^{-1/2}*K_{1/2}^{1/2}
 =\frac1\pi\bcPV\int_0^1
      \frac{K_1^{t-1/2}}{\sqrt{t(1-t)}}\,dt+\frac1\pi.
\end{equation}
The zero Fourier coefficient of the convolution is $1/\pi$, whereas
that of the principal-value mixture is zero. Thus even a single mode
disproves the unrestricted, uncorrected beta formula. This is an exact
counterexample, independent of floating-point calculations.

For a genuinely three-factor example, let the twists be
$(-1/2,1/2,3/2)$ and all three orders be $1/3$. The total order is one;
the crossed modes are $0$ and $-1$. Both have
$Q_n=2\,3^{-1/3}$ and $\sin(\pi\sigma_n)=\sqrt3/2$. Therefore the
correction is
\begin{equation}\label{bc:eq:three-example}
 \frac{3^{1/6}}{2\pi}\left(1+\bce^{-2\pi i x}\right).
\end{equation}
Repeated twists cause no difficulty: identical vertices may be grouped
by adding their orders, directly from the beta integral or the Fourier
multipliers.

\section{All balanced order jets and their chamber rank}
\label{bc:sec:jets}

Fix positive integers $r_j$ with $\sum_jr_j=m$. Choose a direction
$\bm v=(v_1,\ldots,v_d)$ with $\sum_jv_j=0$, and put
$\alpha_j=r_j+uv_j$. For sufficiently small $u$ the hypotheses of
Theorem~\ref{bc:thm:periodic} hold. At a crossed mode define
\begin{align}
 R_n&=\sum_{n+\theta_j<0}r_j,&
 V_n&=\sum_{n+\theta_j<0}v_j,\notag\\
 L_n&=-\sum_jv_j\log|n+\theta_j|,&
 Q_{n,0}&=\prod_j|n+\theta_j|^{-r_j}.
 \label{bc:eq:jet-data}
\end{align}
For a nonnegative integer $q$, call $C_q$ the $q$th $u$ derivative at zero
of the correction in \eqref{bc:eq:periodic-main}.

\begin{theorem}[Every balanced derivative, without a recursion]
\label{bc:thm:jets}
For every $q\ge0$,
\begin{equation}\label{bc:eq:jet-main}
 C_q(x)=\frac1{2(2\pi i)^m}\sum_{n\in\bcC}
 (-1)^{R_n}Q_{n,0}\bce^{2\pi i n x}
 \left\{(L_n+i\pi V_n)^q-(L_n-i\pi V_n)^q\right\}.
\end{equation}
In particular $C_0=0$, but
\begin{equation}\label{bc:eq:first-jet}
 C_1(x)=\frac{i\pi}{(2\pi i)^m}\sum_{n\in\bcC}
             (-1)^{R_n}Q_{n,0}V_n\bce^{2\pi i n x}.
\end{equation}
\end{theorem}
\begin{proof}
At each crossed mode the varying factor is
\[
 i\sin(\pi(R_n+uV_n))\,Q_{n,0}\bce^{uL_n}
 =\frac{(-1)^{R_n}Q_{n,0}}2
     \left(\bce^{u(L_n+i\pi V_n)}
                  -\bce^{u(L_n-i\pi V_n)}\right).
\]
Differentiate this finite sum. Holomorphic dependence in the preceding
theorem justifies differentiating the convolution and mixture as well.
\end{proof}

For several balanced directions $\bm v^{(1)},\ldots,\bm v^{(p)}$, replace
$\alpha_j$ by $r_j+\sum_\ell u_\ell v_j^{(\ell)}$. Define $L_{n,\ell}$ and
$V_{n,\ell}$ as in \eqref{bc:eq:jet-data}. The mixed derivative of
multi-order $\bm q$ has the same formula with the brace replaced by
\begin{equation}\label{bc:eq:multi-jet}
 \prod_{\ell=1}^p(L_{n,\ell}+i\pi V_{n,\ell})^{q_\ell}
 -\prod_{\ell=1}^p(L_{n,\ell}-i\pi V_{n,\ell})^{q_\ell}.
\end{equation}
Thus arbitrarily high mixed derivatives require only finite
polynomials in logarithms and $\pi$ on the correction side.

\subsection{Harmonic-number weights on the simplex}

Let $\bcY_q$ denote the complete exponential Bell polynomial, normalized by
\[
 \exp\left(\sum_{k\ge1}\ell_k\frac{u^k}{k!}\right)
          =\sum_{q\ge0}\bcY_q(\ell_1,\ldots,\ell_q)\frac{u^q}{q!}.
\]
For $N\ge0$ and $k\ge1$, put
$H_N^{(k)}=\sum_{j=1}^N j^{-k}$, with $H_0^{(k)}=0$ and $H_N=H_N^{(1)}$.
The simplex density at the integer tuple is
\[
 B_0(\bm t)=\frac{\Gamma(m)}{\prod_j\Gamma(r_j)}
                             \prod_jt_j^{r_j-1}.
\]
Its $q$th balanced derivative is
\begin{equation}\label{bc:eq:bell-weight}
 B_0(\bm t)\bcY_q(\ell_1,\ldots,\ell_q),
\end{equation}
where
\begin{align}
 \ell_1&=\sum_jv_j(\log t_j-H_{r_j-1}),
       \label{bc:eq:ell-one}\\
 \ell_k&=(-1)^{k+1}(k-1)!
           \sum_jv_j^k\bigl(\zeta(k)-H_{r_j-1}^{(k)}\bigr),
                      \qquad k\ge2.
       \label{bc:eq:ell-higher}
\end{align}
To prove these expressions, differentiate the logarithm of the normalized
Dirichlet density. Its first derivative is
$\sum_jv_j(\log t_j-\psi(r_j))$. The Euler constant cancels since
$\sum v_j=0$ and $\psi(r_j)=H_{r_j-1}-\gamma$. Higher derivatives are
$-\sum_jv_j^k\psi^{(k-1)}(r_j)$. The convergent polygamma series gives
\eqref{bc:eq:ell-higher}. Exponentiation proves
\eqref{bc:eq:bell-weight}.

Consequently the derivative of the mixture is the affine finite-part
integral of $K_m^{\sum t_j\theta_j}$ against
\eqref{bc:eq:bell-weight}. These formulas evaluate whole families of
logarithmic simplex integrals. They are not merely derivatives of an
identity at an isolated integer point: the required analytic family has
already been established.

\begin{example}[A zero correction with a nonzero derivative]
For $\theta_1=-1/2$, $\theta_2=1/2$ and
$(\alpha_1,\alpha_2)=(1+u,1-u)$, the zero-mode coefficient of the
convolution is $\pi^{-2}\bce^{i\pi u}$. The corresponding finite-part
coefficient is $\pi^{-2}\cos(\pi u)$. Their difference is
\begin{equation}\label{bc:eq:jet-example}
 \frac{i\sin(\pi u)}{\pi^2}.
\end{equation}
It is zero at $u=0$, but its first and third derivatives there are
$i/\pi$ and $-i\pi$. The second derivative is zero. An integer-order
partial-fraction formula alone cannot supply these order derivatives.
\end{example}

\subsection{Exactly which directions have no correction?}

Let $c$ be the number of occupied open unit intervals $(k,k+1)$ among the
twists, and let $I_1,\ldots,I_c$ be their groups, ordered from left to
right. Empty intervals are not counted.

\begin{theorem}[Chamberwise kernel and rank]\label{bc:thm:rank}
On the balanced tangent space $\{\bm v:\sum_jv_j=0\}$, the linear map
$\bm v\mapsto C_1$ has rank $c-1$. Its kernel is exactly
\begin{equation}\label{bc:eq:chamber-kernel}
 \left\{\bm v:\ \sum_{j\in I_h}v_j=0\text{ for every }h=1,\ldots,c\right\},
\end{equation}
and has dimension $d-c$. A direction belongs to this kernel if and only
if \emph{every} $C_q$ vanishes. In that case the full correction vanishes
identically as a function of $u$ near zero.
\end{theorem}
\begin{proof}
Because the exponential functions at distinct integer frequencies are
linearly independent and $Q_{n,0}\ne0$, formula
\eqref{bc:eq:first-jet} vanishes if and only if $V_n=0$ at each crossed
mode. For $n=-k$, the negative group is $\{j:\theta_j<k\}$. It is a union
of complete occupied chambers. As the integer $k$ runs between the
smallest and largest twists, these groups realize every nonempty proper
prefix $I_1\cup\cdots\cup I_h$, $1\le h<c$. Gaps between occupied
chambers repeat a prefix but do not create an additional one.

Writing $s_h=\sum_{j\in I_h}v_j$, the distinct conditions are
$s_1=0$, $s_1+s_2=0$, \ldots, $s_1+\cdots+s_{c-1}=0$. Together with
$\sum_hs_h=0$, these are equivalent to all $s_h=0$. The map from $\bm v$
to the chamber sums is onto the $(c-1)$-dimensional zero-sum space:
choose one representative index per chamber. Hence the rank is $c-1$
and the kernel dimension is $d-c$. If all $V_n$ are zero, the sine
factor $\sin(\pi(R_n+uV_n))$ is identically zero; thus every derivative
vanishes. The converse already follows from $q=1$.
\end{proof}

For example, twists $(1/5,7/10,6/5,18/5)$ occupy three chambers. The rank
is two and the kernel is generated by $(1,-1,0,0)$. This rank is purely
spectral and exact. It gives no information about linear independence of
special values obtained by evaluating the kernels at a point.

\section{Centered logarithmic moments and finite even-zeta reductions}
\label{bc:sec:moments}

The same crossing calculation produces ordinary real integrals with
arbitrary logarithmic order. Fix $A,B>0$ and define
\begin{equation}\label{bc:eq:centered-data}
 U(t)=\log\frac{At}{B(1-t)},\qquad
 D(t)=B-(A+B)t,\qquad t_0=\frac{B}{A+B}.
\end{equation}
Both $U$ and $D$ vanish simply at $t_0$. We initially allow positive real
$r,s$ such that $r+s=m\in\bcN$. The stronger polynomial conclusions will
assume that $r,s$ themselves are positive integers.

\begin{theorem}[The centered generator and all ordinary moments]
\label{bc:thm:moments}
For $-r<\bcRe u<s$,
\begin{equation}\label{bc:eq:centered-generator}
 I_{r,s}(u):=\bcFP\int_0^1
 \frac{t^{r-1}(1-t)^{s-1}\bce^{uU(t)}}{D(t)^m}\,dt
 =\frac{\Gamma(r+u)\Gamma(s-u)\cos(\pi(r+u))}
        {(m-1)!A^rB^s}.
\end{equation}
The finite part is in the affine variable $D$. If $r,s\in\bcN$, put
\begin{equation}\label{bc:eq:P}
 P_{r,s}(u)=\prod_{j=1}^{r-1}(j+u)
                \prod_{j=1}^{s-1}(j-u),\qquad
 G_{r,s}(u)=\pi u\cot(\pi u)P_{r,s}(u).
\end{equation}
Empty products equal one. Then
\begin{equation}\label{bc:eq:cot-generator}
 I_{r,s}(u)=\frac{(-1)^r}{(m-1)!A^rB^s}G_{r,s}(u).
\end{equation}
For every integer $q\ge m$, the following integral is ordinary and
absolutely convergent:
\begin{align}
 M_q^{r,s}(A,B)
 &:={}\int_0^1\frac{t^{r-1}(1-t)^{s-1}U(t)^q}{D(t)^m}\,dt\notag\\
 &=\frac{2(-1)^{r+1}q!}{(m-1)!A^rB^s}
       \sum_{k=1}^{\lfloor q/2\rfloor}
           \zeta(2k)[u^{q-2k}]P_{r,s}(u).
 \label{bc:eq:all-moments}
\end{align}
\end{theorem}
\begin{proof}
Since
$\bce^{uU(t)}=(A/B)^u t^u(1-t)^{-u}$, apply
\eqref{bc:eq:two-beta} with $\alpha=r+u$ and $\beta=s-u$.
The factors $A^{-u}B^u$ on its right cancel the centering factor
$(A/B)^u$. This proves \eqref{bc:eq:centered-generator}. Endpoint
integrability on compact subsets of the indicated strip, and Taylor
subtraction at $t_0$, show that both sides are holomorphic there.

For integer $r,s$, the Gamma recurrence and reflection identity give
\[
 \Gamma(r+u)\Gamma(s-u)
  =P_{r,s}(u)\Gamma(1+u)\Gamma(1-u)
  =P_{r,s}(u)\frac{\pi u}{\sin\pi u}.
\]
Because $\cos(\pi(r+u))=(-1)^r\cos\pi u$, this proves
\eqref{bc:eq:cot-generator}. The apparent poles at the nonzero integers
inside the strip are canceled by $P_{r,s}$; the value at zero is
removable.

Near $t_0$, $U(t)^q/D(t)^m=O(|t-t_0|^{q-m})$. For $q\ge m$ this is
bounded, and the endpoint logarithms are integrable because $r,s>0$.
Thus the finite-part derivative $I_{r,s}^{(q)}(0)$ equals the ordinary
integral in \eqref{bc:eq:all-moments}. Finally, the logarithmic derivative
of the product for $\sin\pi u$, or the partial fractions of cotangent,
gives the convergent expansion
\begin{equation}\label{bc:eq:cot-series}
 \pi u\cot\pi u=1-2\sum_{k\ge1}\zeta(2k)u^{2k},\qquad |u|<1.
\end{equation}
For example, the partial fractions are
$\pi u\cot\pi u=1-2\sum_{n\ge1}u^2/(n^2-u^2)$, and expanding the
absolutely convergent right side proves \eqref{bc:eq:cot-series}.
Since $\deg P_{r,s}=m-2<q$, the initial polynomial contributes no
$u^q$ coefficient. Coefficient extraction proves the stated formula.
\end{proof}

The endpoint asymmetry has disappeared from the answer except for the
factor $A^{-r}B^{-s}$. The centered logarithm is essential for this
simplification. An uncentered logarithm gives derivatives of an
additional exponential in $u$.

\subsection{Harmonic coefficients and parity}

Every coefficient of $P_{r,s}$ is explicit in generalized harmonic
numbers. Namely,
\begin{equation}\label{bc:eq:P-bell}
 [u^j]P_{r,s}(u)=\frac{(r-1)!(s-1)!}{j!}
               \bcY_j(p_1,\ldots,p_j),
\end{equation}
where
\begin{equation}\label{bc:eq:p-harmonic}
 p_k=(k-1)!\left((-1)^{k-1}H_{r-1}^{(k)}-H_{s-1}^{(k)}\right).
\end{equation}
To verify this, take the logarithm of the two finite products in
\eqref{bc:eq:P} and expand each $\log(1\pm u/j)$. This is an identity
of finite Taylor coefficients near zero. Although the Bell expressions
are written for every $j$, the coefficients vanish for $j>m-2$.

For $r=s$, $P_{r,r}(u)=\prod_{j=1}^{r-1}(j^2-u^2)$ is even. Therefore
all odd centered moments vanish. The same fact follows directly from
$t\mapsto1-t$ when $A=B$, and from the logistic coordinate used below
for arbitrary positive $A,B$.

Some explicit ordinary evaluations, with $A=B=1$ and
$U=\log(t/(1-t))$, are
\begin{align}
 \int_0^1\frac{U^{2k}}{(1-2t)^2}\,dt
     &=2(2k)!\zeta(2k),\qquad k\ge1,
                      \label{bc:eq:moment-11}\\
 \int_0^1\frac{(1-t)U^3}{(1-2t)^3}\,dt&=-\pi^2,
                      \label{bc:eq:moment-12a}\\
 \int_0^1\frac{(1-t)U^4}{(1-2t)^3}\,dt&=\frac{4\pi^4}{15},
                      \label{bc:eq:moment-12b}\\
 \int_0^1\frac{t(1-t)U^4}{(1-2t)^4}\,dt
       &=8\bigl(\zeta(2)-\zeta(4)\bigr),
                      \label{bc:eq:moment-22a}\\
 \int_0^1\frac{t(1-t)U^6}{(1-2t)^4}\,dt
       &=240\bigl(\zeta(4)-\zeta(6)\bigr).
                      \label{bc:eq:moment-22b}
\end{align}
For instance $P_{1,2}=1-u$ and $P_{2,2}=1-u^2$, so each evaluation is
an immediate finite coefficient extraction. None uses a numerical
integer-relation search.

\subsection{An ordinary subtracted generating integral}

Let $T_NF(u)=\sum_{q=0}^NF^{(q)}(0)u^q/q!$ be the Taylor polynomial
at zero. Write $c_{r,s}=(-1)^r/((m-1)!A^rB^s)$. The finite-part generator
can be converted into an ordinary one:
\begin{corollary}[Subtraction of the full local obstruction]
\label{bc:cor:ordinary-generator}
For $r,s\in\bcN$ and $-r<\bcRe u<s$,
\begin{align}
 &\int_0^1\frac{t^{r-1}(1-t)^{s-1}}{D(t)^m}
     \left(\bce^{uU(t)}-\sum_{q=0}^{m-1}
                        \frac{u^qU(t)^q}{q!}\right)dt\notag\\
 &\hspace{35mm}=c_{r,s}\bigl(G_{r,s}(u)-T_{m-1}G_{r,s}(u)\bigr).
 \label{bc:eq:ordinary-generator}
\end{align}
The integral is ordinary, locally uniformly convergent in the strip.
\end{corollary}
\begin{proof}
At $t_0$ the numerator is $O(U^m)$, uniformly on compact $u$ sets.
At the endpoints, the exponential term is bounded by powers with
positive integrability exponents on every compact substrip, and the
subtractions have only logarithmic growth. Taylor subtraction in
\eqref{bc:eq:centered-generator} therefore agrees with the ordinary
integral on the left. Equation \eqref{bc:eq:cot-generator} gives the
right side.
\end{proof}

\section{Classical-polylogarithm primitives and inverse-logarithm integrals}
\label{bc:sec:polylog}

The generating function in \eqref{bc:eq:cot-generator} has elementary
meromorphic factors, but its antiderivatives naturally introduce
polylogarithms. The following finite formula avoids a recursive chain of
unevaluated integrals.

\begin{proposition}[A finite primitive for every polynomial cotangent]
\label{bc:prop:polyprimitive}
Let $R$ be a polynomial of degree $d$. On any sufficiently small simply
connected domain avoiding the integers, choose consistent local
branches of $\bcLi_j(\bce^{2\pi i u})$. Then
\begin{equation}\label{bc:eq:polyprimitive}
 \mathcal A_R(u)=
 -i\pi\int R(u)\,du
 -\sum_{j=0}^{d}\frac{(-1)^jR^{(j)}(u)}{(2\pi i)^j}
                       \bcLi_{j+1}(\bce^{2\pi i u})
\end{equation}
is an antiderivative of $\pi R(u)\cot\pi u$.
\end{proposition}
\begin{proof}
Put $z=\bce^{2\pi i u}$ and $c=2\pi i$. Since
$\bcLi_0(z)=z/(1-z)$, the elementary cotangent formula gives
\[
 \cot\pi u=-i\bigl(1+2\bcLi_0(z)\bigr).
\]
The differential identities
$\partial_u\bcLi_{j+1}(z)=c\bcLi_j(z)$ follow from the power series
in $|z|<1$ and then by analytic continuation on the chosen local
branches. Differentiating the finite sum in
\eqref{bc:eq:polyprimitive}, adjacent terms cancel. The only surviving
polylogarithmic term is $-cR(u)\bcLi_0(z)$, while differentiating the
polynomial primitive contributes $-i\pi R(u)$. Their sum is
$\pi R(u)\cot\pi u$.
\end{proof}

This is a local primitive, not a globally single-valued function on the
punctured plane. In particular the logarithm in $\bcLi_1(z)=-\log(1-z)$
requires a branch. For the definite real integrals below, analytic
continuation from the upper half $u$-plane gives an unambiguous boundary
value on $0<u<1$.

\begin{corollary}[Normalized real primitive]
\label{bc:cor:real-primitive}
Suppose $R$ has real coefficients and $R(0)=0$. For $0<v<1$,
\begin{align}
 \int_0^v\pi R(u)\cot\pi u\,du
 ={}&-\bcRe\sum_{j=0}^{d}
       \frac{(-1)^jR^{(j)}(v)}{(2\pi i)^j}
                           \bcLi_{j+1}(\bce^{2\pi i v})\notag\\
 &+\bcRe\sum_{j=1}^{d}
       \frac{(-1)^jR^{(j)}(0)}{(2\pi i)^j}\zeta(j+1).
 \label{bc:eq:real-primitive}
\end{align}
If additionally $v=1/2$, this becomes the finite odd-zeta expression
\begin{align}
 \int_0^{1/2}\pi R(u)\cot\pi u\,du
 ={}&R(1/2)\log2\notag\\
 &+\sum_{\ell=1}^{\lfloor d/2\rfloor}
   \frac{(-1)^\ell}{(2\pi)^{2\ell}}
   \left((1-2^{-2\ell})R^{(2\ell)}(1/2)+R^{(2\ell)}(0)\right)
                           \zeta(2\ell+1).
 \label{bc:eq:half-primitive}
\end{align}
\end{corollary}
\begin{proof}
The real part of $-i\pi\int R$ is zero along the real segment. As
$u\downarrow0$, the $j=0$ term tends to zero since
$R(u)=O(u)$ and $\bcLi_1(\bce^{2\pi i u})=O(1+|\log u|)$.
For $j\ge1$, the polylogarithm tends to $\zeta(j+1)$. Subtract the
endpoint value in Proposition~\ref{bc:prop:polyprimitive}. This proves
\eqref{bc:eq:real-primitive}. At $v=1/2$, use
$\bcLi_1(-1)=-\log2$ and
$\bcLi_{j+1}(-1)=-(1-2^{-j})\zeta(j+1)$ for $j\ge1$.
The odd-$j$ terms have zero real part, giving
\eqref{bc:eq:half-primitive}.
\end{proof}

For $R(u)=uP_{r,s}(u)$, the degree is $m-1$ and
$\pi R\cot\pi u=G_{r,s}$. Thus polylogarithms through order $m$ suffice
for the primitive of every centered generator. In the symmetric example
$r=s=2$, $R(u)=u-u^3$, and
\begin{equation}\label{bc:eq:primitive-22}
 \int_0^{1/2}\pi(u-u^3)\cot\pi u\,du
      =\frac38\log2+\frac{9\zeta(3)}{16\pi^2}.
\end{equation}
The appearance of odd zeta values in a primitive is compatible with the
even-zeta Taylor coefficients of the original generator: integration
and evaluation at a nonzero endpoint are different operations.

\subsection{An all-order ordinary inverse-logarithm family}

Integrating \eqref{bc:eq:ordinary-generator} in the parameter, rather than
in $t$, gives the following exact identity.

\begin{theorem}[Convergent inverse-logarithm integrals]
\label{bc:thm:inverse-log}
Let $r,s\in\bcN$, $m=r+s$, and $0<v<1$. Then
\begin{align}
 &\int_0^1\frac{t^{r-1}(1-t)^{s-1}}{D(t)^m}
 \frac{\displaystyle\bce^{vU(t)}
              -\sum_{j=0}^{m}\frac{(vU(t))^j}{j!}}{U(t)}\,dt\notag\\
 &\quad=c_{r,s}\left\{\int_0^vG_{r,s}(u)\,du
       -\sum_{q=0}^{m-1}
          \frac{G_{r,s}^{(q)}(0)v^{q+1}}{(q+1)!}\right\}.
 \label{bc:eq:inverse-log}
\end{align}
The left side is an ordinary absolutely convergent integral with a
removable expression at $t=t_0$. Its right side is a finite combination
of classical polylogarithms and polynomial terms by
\eqref{bc:eq:real-primitive}. At $v=1/2$, only $\log2$, odd zeta values,
even zeta values from the Taylor polynomial, and rational powers of
$\pi$ are needed.
\end{theorem}
\begin{proof}
Integrate the left side of \eqref{bc:eq:ordinary-generator} over
$0\le u\le v$. Uniform domination on this compact parameter segment
justifies Fubini. The integral of the exponential minus its first $m$
Taylor terms is exactly the displayed quotient. Its numerator vanishes
to order $m+1$ in $U$, so division by $U D^m$ leaves a bounded function
at $t_0$. The endpoint convergence follows from $v<s$, automatic here
because $s\ge1$ and $v<1$. Integrating the right side proves the formula.
\end{proof}

Here are two explicit evaluations. In both, $U=\log(t/(1-t))$:
\begin{align}
 &\int_0^1\frac{\bce^{U/2}-1-U/2-U^2/8}{U(1-2t)^2}\,dt
             =\frac{1-\log2}{2},
 \label{bc:eq:inverse-11}\\[3pt]
 &\int_0^1\frac{t(1-t)}{(1-2t)^4}
  \frac{\bce^{U/2}-1-U/2-U^2/8-U^3/48-U^4/384}{U}\,dt\notag\\
 &\hspace{18mm}=
  \frac{\log2}{16}+\frac{3\zeta(3)}{32\pi^2}
                   -\frac{11}{144}+\frac{\pi^2}{432}.
 \label{bc:eq:inverse-22}
\end{align}
For the second identity,
$G_{2,2}(u)=1-(1+\pi^2/3)u^2+O(u^4)$, so the integrated Taylor
polynomial is $11/24-\pi^2/72$ at $v=1/2$. Combine this with
\eqref{bc:eq:primitive-22} and $c_{2,2}=1/6$. This supplies an entirely
symbolic derivation of the constant.

\section{A logistic Mellin transform and finite Hurwitz reductions}
\label{bc:sec:mellin}

The centered logarithm is also a natural global coordinate. Solving
\eqref{bc:eq:centered-data} for $t$ gives
\begin{equation}\label{bc:eq:logistic}
 t=\frac{B\bce^U}{A+B\bce^U},\qquad
 \frac{t^{r-1}(1-t)^{s-1}}{D(t)^m}\,dt
   =\frac1{A^rB^s}\frac{\bce^{rU}}{(1-\bce^U)^m}\,dU.
\end{equation}
This identity holds as an ordinary change of variables away from $U=0$.
For the convergent moments $q\ge m$, it is valid across $U=0$ as well.
On $U<0$ put $v=-U$, and on $U>0$ put $v=U$. The two kernels become
$\bce^{-rv}/(1-\bce^{-v})^m$ and
$(-1)^m\bce^{-sv}/(1-\bce^{-v})^m$.

For an integer $m\ge1$ and $a>0$, define
\begin{align}
 \bcZ_m(z,a)&=\sum_{n=0}^\infty
       \binom{n+m-1}{m-1}(n+a)^{-z},\qquad \bcRe z>m,
                  \label{bc:eq:barnes-series}\\
 \bcJ_{m,a}(z)&=\int_0^\infty
       v^{z-1}\frac{\bce^{-av}}{(1-\bce^{-v})^m}\,dv,
                  \qquad \bcRe z>m.
                  \label{bc:eq:mellin-J}
\end{align}
The function $\bcZ_m$ is the equal-period Barnes-type zeta: the
multiplicity counts the nonnegative integer $m$-tuples of sum $n$.
We use only the explicit finite reduction proved next, not any
unproved reduction of a general multiple zeta function.

Write
\begin{equation}\label{bc:eq:c-polynomial}
 C_{m,a}(X)=\prod_{j=1}^{m-1}(X+j-a)
                  =\sum_{k=0}^{m-1}c_{m,a,k}X^k.
\end{equation}
For $m=1$ this polynomial is one.

\begin{proposition}[Mellin transform, order derivatives, and polygamma]
\label{bc:prop:mellin}
For $a>0$,
\begin{equation}\label{bc:eq:finite-hurwitz}
 \bcZ_m(z,a)=\frac1{(m-1)!}
                 \sum_{k=0}^{m-1}c_{m,a,k}\zeta(z-k,a),
 \qquad
 \bcJ_{m,a}(z)=\Gamma(z)\bcZ_m(z,a).
\end{equation}
The first equality gives meromorphic continuation to the whole
$z$-plane. For every $q\ge0$ and $\bcRe z>m$,
\begin{equation}\label{bc:eq:mellin-log}
 \int_0^\infty v^{z-1}(\log v)^q
              \frac{\bce^{-av}}{(1-\bce^{-v})^m}\,dv
 =\partial_z^q\left(\Gamma(z)\bcZ_m(z,a)\right).
\end{equation}
At an integer $p\ge m$, the non-logarithmic case is the finite
polygamma expression
\begin{equation}\label{bc:eq:polygamma-Mellin}
 \bcJ_{m,a}(p+1)=\frac{p!}{(m-1)!}
       \sum_{k=0}^{m-1}c_{m,a,k}
            \frac{(-1)^{p+1-k}}{(p-k)!}\psi^{(p-k)}(a).
\end{equation}
\end{proposition}
\begin{proof}
Since $C_{m,a}(n+a)=\prod_{j=1}^{m-1}(n+j)$, expanding the polynomial
in \eqref{bc:eq:barnes-series} gives the finite Hurwitz sum in its
absolutely convergent half-plane. Next expand
$(1-\bce^{-v})^{-m}=\sum_{n\ge0}\binom{n+m-1}{m-1}\bce^{-nv}$.
For real $z>m$, Tonelli and the Gamma integral prove the second identity;
absolute convergence and holomorphy extend it to $\bcRe z>m$.
Near zero the integrand is $O(v^{\bcRe z-m-1})$, while at infinity
it decays exponentially. Extra powers of $\log v$ are integrable on
compact subsets of the half-plane, proving
\eqref{bc:eq:mellin-log} by differentiation. Finally, the convergent
polygamma series
$\psi^{(h)}(a)=(-1)^{h+1}h!\zeta(h+1,a)$ gives
\eqref{bc:eq:polygamma-Mellin}.
\end{proof}

This proof also makes clear which parts of the construction are
classical: the finite Hurwitz reduction is polynomial multiplicity
expansion, and the Mellin kernel is its elementary Gamma transform.
Their purpose here is to give a second derivation of the moment family
and a precise Stieltjes finite-part continuation.

\subsection{A second evaluation of all centered moments}

For $r,s>0$, $r+s=m\in\bcN$ and $q\ge m$ integral,
\eqref{bc:eq:logistic} gives
\begin{equation}\label{bc:eq:moment-second}
 M_q^{r,s}(A,B)=\frac1{A^rB^s}
    \left((-1)^q\bcJ_{m,r}(q+1)+(-1)^m\bcJ_{m,s}(q+1)\right).
\end{equation}
For integer $r,s$ this is a finite combination of generalized harmonic
tails. Equating \eqref{bc:eq:moment-second} with
\eqref{bc:eq:all-moments} proves the finite identity
\begin{align}
 &(-1)^q\sum_{j=0}^{m-1}c_{m,r,j}\zeta(q+1-j,r)
  +(-1)^m\sum_{j=0}^{m-1}c_{m,s,j}\zeta(q+1-j,s)\notag\\
 &\qquad=2(-1)^{r+1}
     \sum_{k=1}^{\lfloor q/2\rfloor}
                   \zeta(2k)[u^{q-2k}]P_{r,s}(u).
 \label{bc:eq:hurwitz-cancellation}
\end{align}
In particular, replacing
$\zeta(h,r)=\zeta(h)-H_{r-1}^{(h)}$ for the integer shifts shows the
exact cancellation of the unwanted odd-zeta and rational harmonic
tail terms. The two analytic proofs establish this at all orders; the
216 symbolic comparisons supplied with the article are independent
finite consistency checks, not the reason it holds for all orders.

There is also an ordinary fractional-order extension on either half
of the $U$-line. For example, for $\bcRe z>m$,
\begin{align}
 \int_0^{t_0}\frac{t^{r-1}(1-t)^{s-1}[-U(t)]^{z-1}}
                        {D(t)^m}\,dt
       &=A^{-r}B^{-s}\Gamma(z)\bcZ_m(z,r),\notag\\
 \int_{t_0}^{1}\frac{t^{r-1}(1-t)^{s-1}U(t)^{z-1}}
                        {D(t)^m}\,dt
       &=(-1)^mA^{-r}B^{-s}\Gamma(z)\bcZ_m(z,s).
 \label{bc:eq:half-Mellin}
\end{align}
The powers of positive $\pm U$ use real logarithms. Their order
derivatives insert $\log(-U)$ or $\log U$; no branch of a negative
real power is hidden in these formulas.

\section{Stieltjes coefficients, finite parts, and their primitives}
\label{bc:sec:stieltjes}

The Laurent expansion convention for generalized Stieltjes constants is
\begin{equation}\label{bc:eq:stieltjes-convention}
 \zeta(1+\varepsilon,a)=\frac1\varepsilon+
     \sum_{q=0}^\infty\frac{(-1)^q\gamma_q(a)}{q!}\varepsilon^q.
\end{equation}
In particular $\gamma_0(a)=-\psi(a)$. All $a$ parameters in this section
are positive real numbers; differentiation may equivalently be carried
out on a complex neighborhood avoiding the nonpositive integers.

For $1\le j\le m$ set
\begin{equation}\label{bc:eq:rho}
 \rho_{m,j}(a)=\frac{c_{m,a,j-1}}{(m-1)!}.
\end{equation}
The coefficient may be zero, in which case the apparent pole is
removable. Define $E_{m,j,q}(a)$ by
\begin{equation}\label{bc:eq:Z-Laurent}
 \bcZ_m(j+\varepsilon,a)=\frac{\rho_{m,j}(a)}\varepsilon
                         +\sum_{q\ge0}E_{m,j,q}(a)\varepsilon^q.
\end{equation}
The finite Hurwitz reduction gives the explicit formula
\begin{equation}\label{bc:eq:E}
 E_{m,j,q}(a)=\frac{1}{(m-1)!q!}
 \left\{c_{m,a,j-1}(-1)^q\gamma_q(a)
       +\sum_{\substack{0\le k\le m-1\\k\ne j-1}}
               c_{m,a,k}\zeta^{(q)}(j-k,a)\right\}.
\end{equation}
The superscript on $\zeta$ here always means differentiation in its
first, order variable. Formula \eqref{bc:eq:E} includes the residue-zero
cases without a limiting convention.

\subsection{The unnormalized Mellin finite part has an extra term}

We now fix the cutoff variable to be $v$ in \eqref{bc:eq:mellin-J}, with
reference split at $v=1$. The notation
\[
 \bcFP\int_0^\infty v^{j-1}(\log v)^q
                   \frac{\bce^{-av}}{(1-\bce^{-v})^m}\,dv
\]
means the Hadamard finite part obtained by subtracting all divergent
terms of the small-$v$ expansion and discarding precisely their
divergent cutoff powers and logarithms. The normalization is not
inferred from the earlier $t$-coordinate finite part.

\begin{theorem}[Every logarithmic finite part at a Mellin pole]
\label{bc:thm:mellin-fp}
For $a>0$, $1\le j\le m$ and $q\ge0$,
\begin{align}
 &\bcFP\int_0^\infty v^{j-1}(\log v)^q
                  \frac{\bce^{-av}}{(1-\bce^{-v})^m}\,dv\notag\\
 &\quad=q!\left\{
     \rho_{m,j}(a)\frac{\Gamma^{(q+1)}(j)}{(q+1)!}
     +\sum_{\ell=0}^q
         \frac{\Gamma^{(\ell)}(j)}{\ell!}
                     E_{m,j,q-\ell}(a)\right\}.
 \label{bc:eq:Mellin-FP}
\end{align}
In particular, the residue contribution involving
$\Gamma^{(q+1)}(j)$ cannot generally be omitted.
\end{theorem}
\begin{proof}
Let
\[
 g(v)=\bce^{-av}\left(\frac{v}{1-\bce^{-v}}\right)^m
                      =\sum_{h\ge0}b_hv^h
\]
near zero. This is analytic with $b_0=1$. Subtract its Taylor polynomial
through $h=m-j$ on $0<v<1$. The remainder of
$v^{j-m-1}g(v)(\log v)^q$ is integrable. Each subtracted monomial with
$h<m-j$ has finite-part integral
\begin{equation}\label{bc:eq:monomial-fp}
 \bcFP\int_0^1v^{j-m+h-1}(\log v)^q\,dv
   =\frac{(-1)^qq!}{(j-m+h)^{q+1}},
\end{equation}
whereas the monomial with $h=m-j$ has finite part zero. This follows
by direct integration from a cutoff $\delta$ to one; in the last case
the antiderivative is $(\log v)^{q+1}/(q+1)$.

Equivalently, insert $v^\varepsilon$ before integration. After the same
subtractions the meromorphic continuation is obtained by adding
$b_h/(\varepsilon+j-m+h)$. Its $q$th derivative in $\varepsilon$
inserts $(\log v)^q$. The constant Laurent coefficient of that derivative
is exactly the finite part calculated in \eqref{bc:eq:monomial-fp}.
Consequently our finite part is
$q![\varepsilon^q]\bcJ_{m,a}(j+\varepsilon)$, where the bracket means
the regular Laurent coefficient. Multiply
\eqref{bc:eq:Z-Laurent} by the Taylor expansion of
$\Gamma(j+\varepsilon)$. Its $\varepsilon^q$ coefficient is
\[
 \rho_{m,j}\frac{\Gamma^{(q+1)}(j)}{(q+1)!}
 +\sum_{\ell=0}^q\frac{\Gamma^{(\ell)}(j)}{\ell!}E_{m,j,q-\ell}.
\]
This proves \eqref{bc:eq:Mellin-FP}, including any removable-pole case.
\end{proof}

The Gamma derivatives in this theorem are finite Bell expressions in
polygamma values:
\begin{equation}\label{bc:eq:Gamma-Bell}
 \Gamma^{(q)}(j)=\Gamma(j)
          \bcY_q\bigl(\psi(j),\psi'(j),\ldots,\psi^{(q-1)}(j)\bigr).
\end{equation}
At positive integer $j$ these reduce to $\gamma$, generalized harmonic
numbers, and ordinary zeta values. Thus \eqref{bc:eq:Mellin-FP} is an
explicit finite expression in the functions requested, at every
logarithmic order.

\begin{example}[A normalization test that detects the missing residue]
For $m=2$ and $a=1$,
$\bcZ_2(z,1)=\zeta(z-1)$ and $\rho_{2,2}=1$. At $j=2$,
$E_{2,2,0}=\gamma$ and $\Gamma'(2)=1-\gamma$. Therefore
\begin{equation}\label{bc:eq:Mellin-test-zero}
 \bcFP\int_0^\infty
           \frac{v\bce^{-v}}{(1-\bce^{-v})^2}\,dv=1,
\end{equation}
not $\gamma$. Directly,
\[
 \int_\delta^\infty\frac{v\bce^{-v}}{(1-\bce^{-v})^2}\,dv
  =\frac{\delta}{\bce^\delta-1}-\log(1-\bce^{-\delta})
  =1-\log\delta+o(1).
\]
For the first logarithmic moment, using
$\Gamma''(2)=\zeta(2)-2\gamma+\gamma^2$, the same theorem gives
\begin{equation}\label{bc:eq:Mellin-test-one}
 \bcFP\int_0^\infty
         \frac{v\log v\,\bce^{-v}}{(1-\bce^{-v})^2}\,dv
       =\frac{\zeta(2)-\gamma^2}{2}-\gamma_1.
\end{equation}
Here $\gamma_1=\gamma_1(1)$ in convention
\eqref{bc:eq:stieltjes-convention}.
\end{example}

For arbitrary $a>0$, the same order-two calculation gives the exact
polygamma finite parts
\begin{align}
 \bcFP\int_0^\infty\frac{v\bce^{-av}}{(1-\bce^{-v})^2}\,dv
   &=1-\gamma-\psi(a)+(1-a)\psi'(a),
       \label{bc:eq:fp-polygamma-two}\\
 \bcFP\int_0^\infty\frac{\bce^{-av}}{(1-\bce^{-v})^2}\,dv
   &=(a-1)(\psi(a)+\gamma)+\frac12-a.
       \label{bc:eq:fp-polygamma-one}
\end{align}
They follow from $C_{2,a}(X)=X+1-a$,
$\zeta(0,a)=1/2-a$, and $\zeta(2,a)=\psi'(a)$.
Both are in the declared $v$-cutoff normalization.

\subsection{Parameter derivatives and exact Stieltjes antiderivatives}

The multiplicity in \eqref{bc:eq:barnes-series} does not depend on $a$.
Thus the series gives, initially in its convergent half-plane,
\begin{equation}\label{bc:eq:a-ladder}
 \partial_a\bcZ_m(z,a)=-z\bcZ_m(z+1,a),\qquad
 \partial_a^p\bcZ_m(z,a)=(-1)^p(z)_p\bcZ_m(z+p,a).
\end{equation}
Meromorphic continuation proves these identities for all $z$. This
short ladder is particularly effective at an order where a pole is
multiplied by a vanishing factor.

\begin{theorem}[A primitive for every Stieltjes coefficient]
\label{bc:thm:primitives}
For all $m\ge1$ and $q\ge0$,
\begin{equation}\label{bc:eq:primitive-j1}
 \frac{d}{da}\left[-\frac{\partial_z^{q+1}\bcZ_m(0,a)}{(q+1)!}\right]
        =E_{m,1,q}(a).
\end{equation}
For $2\le j\le m$, define
\begin{align}
 F_{m,j,q}(a)={}&-
  \sum_{\ell=0}^{q}\frac{(-1)^\ell E_{m,j-1,q-\ell}(a)}{(j-1)^{\ell+1}}
  -\frac{(-1)^{q+1}\rho_{m,j-1}(a)}{(j-1)^{q+2}}.
 \label{bc:eq:primitive-higher}
\end{align}
Then $F_{m,j,q}'(a)=E_{m,j,q}(a)$. Consequently every Laurent
coefficient in \eqref{bc:eq:E} has an explicit finite primitive in
Hurwitz-order derivatives and Stieltjes coefficients. Definite
integrals between positive endpoints are differences of the displayed
primitives.
\end{theorem}
\begin{proof}
Expand \eqref{bc:eq:a-ladder} at $z=\varepsilon$:
\[
 \partial_a\bcZ_m(\varepsilon,a)
   =-\varepsilon\bcZ_m(1+\varepsilon,a)
   =-\rho_{m,1}(a)-\sum_{q\ge0}E_{m,1,q}(a)\varepsilon^{q+1}.
\]
Since $\bcZ_m$ is regular at $z=0$, comparison of the
$\varepsilon^{q+1}$ coefficient proves
\eqref{bc:eq:primitive-j1}. For $j\ge2$, use
\[
 \partial_a\left(-\frac{\bcZ_m(j-1+\varepsilon,a)}{j-1+\varepsilon}\right)
     =\bcZ_m(j+\varepsilon,a).
\]
The regular $\varepsilon^q$ coefficient on the left is
\eqref{bc:eq:primitive-higher}; its final term comes from the residue at
$j-1$. Coefficient comparison proves the derivative assertion. All
expansions are locally uniform in $a>0$, as follows either from the
finite Hurwitz representation or from analyticity before continuation.
\end{proof}

For $m=1$, $\bcZ_1(z,a)=\zeta(z,a)$ and
$E_{1,1,q}=(-1)^q\gamma_q(a)/q!$. Thus the familiar Stieltjes primitive
appears as the first member of the ladder:
\begin{equation}\label{bc:eq:classical-Stieltjes-primitive}
 \int\gamma_q(a)\,da=
       \frac{(-1)^{q+1}}{q+1}\zeta^{(q+1)}(0,a)+\text{constant}.
\end{equation}
For $m=2$,
\begin{equation}\label{bc:eq:Z2}
 \bcZ_2(z,a)=\zeta(z-1,a)+(1-a)\zeta(z,a),
\end{equation}
so \eqref{bc:eq:primitive-j1} gives the nontrivial weighted example
\begin{align}
 \frac{d}{da}\left[-\frac12\left\{
             \zeta''(-1,a)+(1-a)\zeta''(0,a)\right\}\right]
   ={}&\log\Gamma(a)-\frac12\log(2\pi)
                           -(1-a)\gamma_1(a).
 \label{bc:eq:weighted-Stieltjes-primitive}
\end{align}
We used $\zeta'(0,a)=\log\Gamma(a)-\tfrac12\log(2\pi)$, with the
standard positive-real normalization \cite{dlmf-hurwitz}. More generally,
\begin{equation}\label{bc:eq:weighted-general-q}
 \frac{d}{da}\left[-\frac{
       \zeta^{(q+1)}(-1,a)+(1-a)\zeta^{(q+1)}(0,a)}{(q+1)!}\right]
 =\frac{\zeta^{(q)}(0,a)+(1-a)(-1)^q\gamma_q(a)}{q!}.
\end{equation}
This is an all-order identity, not an extrapolation from its first
few numerical values. Higher $m$ yields polynomially weighted
Stieltjes combinations with a finite list of shifted Hurwitz-order
derivatives.
% INTEGRATION-END

\section{Corrections, compatibility, and integration recommendations}
\label{bc:sec:audit}

\subsection{Targeted audit outcome}

No mathematical error was found in the common-interval mixture or the
integer-order partial-fraction theorem of the inspected section
\cite{repo-mixed}. The present result should be added as a continuation,
not used to replace their statements with an alleged correction. Its
hypotheses are different: arbitrary nonintegral real twists, but positive
integer total order for the finite-part formula. The source's unrestricted
fractional \emph{total} order in a common interval is retained.

Three tempting extensions are rigorously rejected in this article.
First, removing the common-interval condition while using an uncorrected
principal value is contradicted by \eqref{bc:eq:half-counterexample}.
Second, inferring vanishing order derivatives from a vanishing
integer-tuple correction is contradicted by \eqref{bc:eq:jet-example}.
Third, replacing the finite part of $\Gamma(z)\bcZ_m(z,a)$ by Gamma times
the regular part of $\bcZ_m$ omits a residue term; the independent cutoff
calculation in \eqref{bc:eq:Mellin-test-zero} detects the error immediately.
These are statements about candidate extensions and normalization rules.
They are not attributed to the current manuscript without textual evidence.

\subsection{Compatibility with the existing proof chain}

At a common integer chamber, Theorem~\ref{bc:thm:periodic} reduces to the
manuscript's beta theorem specialized to integer total order. At a positive
integer tuple, it becomes an affine finite-part interpretation of the
same convolution already given by the source's finite resolvent
identities. These are consistency checks at the overlap, not separate
period conjectures. The harmonic Bell weights in
\eqref{bc:eq:ell-one}--\eqref{bc:eq:ell-higher} are compatible with the
finite-product and Gamma-coefficient conventions in the harmonic
section \cite{repo-harmonic}; we do not claim to re-prove a previously
unresolved Choi case.

The name ``correction'' is used in two different senses that must stay
separate in integration. The sine polynomial corrects a beta mixture
across integer \emph{twists}. The Gamma-derivative term corrects a
\emph{Laurent coefficient extraction} through a pole. Neither is an
unannounced Dirac distribution in the spatial variable. Furthermore,
the ordinary moment identities and the subtracted generators must not
be relabeled as regularized integrals: their vanishing numerators make
them genuinely convergent.

\subsection{Suggested placement and review status}

The mathematical core is supplied as
\texttt{integration/07-beta-crossing-identities.tex}, with namespace
\texttt{bc:} for all labels and a separate small preamble fragment.
Its natural placement is after
\texttt{07-mixed-twist-identities.tex} in the Hurwitz--Stieltjes material.
The Mellin and primitive subsections can instead be cross-referenced
from the integration/differentiation chapters. A machine-readable status
ledger lists the new statements, their proof locations, overlap with
classical ingredients, and verification coverage.

The supplied proofs are mathematical arguments, not proof-assistant
formalizations. They have not undergone independent peer review. The
replay scripts test finite algebra and selected numerical instances;
they do not certify every continuum statement. The generated package
changes no repository files. It is intended for normal analytic audit
and integration, preserving the source's existing status distinctions.

\section{Further research questions and proposed next steps}
\label{bc:sec:questions}

The following questions are deliberately phrased as open tasks rather
than conclusions inferred from the finite tests.

\begin{enumerate}
\item \textbf{Noninteger total order across chambers.}
For $\sum\alpha_j\notin\bcN$, determine a comparably useful explicit
real-boundary or finite-part mixture. The right-half-plane regularized
Dirichlet formula remains available, but the two boundary distributions
are no longer a real rational finite part plus a single delta derivative.
A successful extension should specify the global branch convention and
recover \eqref{bc:eq:periodic-main} at every integer total order.

\item \textbf{Unbalanced order derivatives.}
Allow $\sum_jv_j\ne0$, so that differentiation changes the total order.
This introduces order derivatives of the intermediate kernel and
logarithmic boundary distributions at the crossing. Determine the exact
replacement for the finite polynomial jet formula, with all local
counterterms specified. The balanced theorem does not settle this case.

\item \textbf{An explicit connection to the manuscript's Stieltjes jets.}
Translate the balanced derivatives into the exact normalization of the
manuscript's twisted Stieltjes families. Identify which Bell
combinations reduce to finite single-Hurwitz expressions and which
require the existing convergent unequal-shift expansions. A finite
spectral correction alone does not prove the desired finite reduction
of the full convolution.

\item \textbf{Moving twists and chamber transitions.}
Differentiate in the twists as a vertex approaches an integer, keeping
track of the divergent individual Fourier coefficient. Determine
canonical subtractions and transition identities that agree from both
neighboring chambers. At a fixed nonintegral twist configuration our
formulas are explicit, but uniform passage through an integral vertex
needs a separate theorem.

\item \textbf{Rational endpoint polylogarithm reductions.}
Formula \eqref{bc:eq:real-primitive} evaluates the centered primitive at
$v=p/q$ through polylogarithms at $q$th roots of unity. Develop a
conductor-aware reduction of the full family, including its Taylor
subtractions, into the canonical cyclotomic basis used elsewhere in the
repository. At $q=2$ this article provides the finite odd-zeta formula.
No claim of minimal depth or period independence is implied.

\item \textbf{Repeated parameter primitives.}
Iterate the construction behind \eqref{bc:eq:inverse-log}. At the $h$th
primitive the natural integrand contains an exponential remainder divided
by $U^h$ and a correspondingly larger subtraction polynomial. Give a
single closed coefficient formula, branch normalization, and exact
root-of-unity endpoint evaluations at every $h$.

\item \textbf{Mellin residues with unequal periods.}
Replace $(1-\bce^{-v})^{-m}$ by a product
$\prod_j(1-\bce^{-\omega_jv})^{-1}$. For commensurable positive periods,
find the precise finite Hurwitz decomposition, including quasipolynomial
multiplicities and every logarithmic finite part. For incommensurable
periods, determine which pieces of the Stieltjes primitive ladder survive
without pretending that the same finite one-variable reduction holds.

\item \textbf{Coordinate-change laws for finite parts.}
Derive an explicit all-order local formula comparing the affine
$t$-finite part with the logarithmic $U$-finite part for nonconvergent
moments. The ordinary $q\ge m$ identities already agree automatically.
The remaining finite discrepancy should be expressed in the finite jet
of the logistic map, rather than absorbed into an unspecified constant.

\item \textbf{Formal proof and exact regression coverage.}
Formalize the scalar boundary-value theorem and chamber-rank statement,
then use them as interfaces for the Fourier and special-function
calculations. In the short term, add exact multiparameter Bell and
higher-pole primitive tests, together with independent interval
quadrature for selected convergent integrals. Neither failed numerical
searches nor a passing finite suite should alter open period statuses.
\end{enumerate}

\appendix
\section{A direct finite proof of the two-point density derivative}
\label{bc:app:finite-density}

The local derivative in the two-vertex case also has a direct algebraic
proof, independent of boundary-value evaluation. Let $\alpha+\beta=m$
and $N=m-1$. For $0<t<1$,
\begin{equation}\label{bc:eq:finite-density-identity}
 \frac{d^{m-1}}{dt^{m-1}}
      \left(t^{\alpha-1}(1-t)^{\beta-1}\right)
 =\left(\prod_{h=1}^{m-1}(\alpha-h)\right)t^{-\beta}(1-t)^{-\alpha}.
\end{equation}
Indeed Leibniz's rule, after extracting
$t^{-\beta}(1-t)^{-\alpha}$, gives
\[
 \sum_{j=0}^{N}\binom Nj
  (\alpha-1)^{\underline j}
  (\beta-1)^{\underline{N-j}}(-1)^{N-j}
  t^{N-j}(1-t)^j.
\]
Here the underlined exponent denotes a falling factorial. Since
$\beta=m-\alpha$, each product of falling factorials and its sign is
$\prod_{h=1}^{N}(\alpha-h)$, independent of $j$. The remaining sum is
$(t+(1-t))^N=1$, proving \eqref{bc:eq:finite-density-identity}.

For the affine pushforward $y=B-(A+B)t$, its density is
\[
 h(y)=\frac{\Gamma(m)}{\Gamma(\alpha)\Gamma(\beta)(A+B)}
 \left(\frac{B-y}{A+B}\right)^{\alpha-1}
 \left(\frac{A+y}{A+B}\right)^{\beta-1}.
\]
Use \eqref{bc:eq:finite-density-identity} and the chain rule at $y=0$.
The powers of $A+B$ cancel. Gamma reflection then yields
\[
 h^{(m-1)}(0)=\frac{\Gamma(m)}{\pi A^\alpha B^\beta}
                                      \sin(\pi\alpha),
\]
which is exactly \eqref{bc:eq:density-derivative}. Apparent singularities
in intermediate Gamma quotients at integer parameters are removable;
the polynomial identity itself has no such restriction. The exact
verification script checks its finite Leibniz expansion through $m=12$.

\section{Reproducible verification and evidence boundaries}
\label{bc:app:verification}

The scripts require Python, SymPy and mpmath. The recorded run used
Python 3.13.5, SymPy 1.14.0 and mpmath 1.3.0. All numerical diagnostics
were run at 65 decimal digits. The following commands can be run from
the package directory:
\begin{verbatim}
python verification/check_exact.py
python verification/check_numeric.py --dps 65
latexmk -pdf -interaction=nonstopmode -halt-on-error article.tex
\end{verbatim}
The two scripts write JSON receipts next to themselves and stop on a
failed assertion. Their numerical tolerances are explicitly recorded
per test. These tolerances are diagnostics, not rigorous error enclosures.

\subsection{Finite exact checks}

\begin{center}
\begin{tabular}{lr}
\toprule
Family & Checks\\
\midrule
Two-point finite density derivative &12\\
Multiplicity polynomial &12\\
Polynomial derivative recurrence &12\\
Individual coefficient derivative recurrence &66\\
Sine/exponential jet versus two powers &14\\
Cosine/exponential jet versus two powers &14\\
Chamber rank over exact rational twists &120\\
Mellin moments versus cotangent coefficients &216\\
Gamma prefactor through a simple pole &11\\
Polylogarithm primitive telescoping &25\\
Explicit phase and normalization constants &3\\
Rejected uncorrected formulas &3\\
\midrule
Total &508\\
\bottomrule
\end{tabular}
\end{center}

The 216 moment comparisons cover $1\le r,s\le6$ and the six orders
$m\le q\le m+5$ for each pair. They treat the zeta values as independent
formal symbols and use exact rational harmonic sums, so accidental
numerical cancellation cannot hide a coefficient error. The 120 rank
checks use rational twists and exact matrix ranks, with a fixed seed.
The proof of Theorem~\ref{bc:thm:rank}, not this finite sample, supplies
the general rank statement.

\subsection{Floating-point diagnostics}

\begin{center}
\begin{tabular}{lr}
\toprule
Comparison & Cases\\
\midrule
Direct affine beta finite parts &6\\
Original-variable ordinary centered moments &27\\
Polylogarithm primitive versus quadrature &16\\
Half-period primitive evaluation &1\\
Ordinary inverse-logarithm integrals &2\\
Mellin integrals and first order derivatives &24\\
Direct subtracted logarithmic Mellin finite parts &12\\
Stieltjes primitive differentiated in the parameter &6\\
\midrule
Total &94\\
\bottomrule
\end{tabular}
\end{center}

The beta finite parts are calculated by Taylor subtraction in the
original affine coordinate, not by fitting values at a small imaginary
regulator. The moment diagnostics integrate in $t$, independently of
the Mellin substitution. The pole diagnostics subtract the small-$v$
series, integrate the remainder, and add the exact monomial finite
parts. These are meaningfully different evaluations from the closed
forms, although floating-point quadrature and finite Taylor arithmetic
still do not produce rigorous intervals.

The supplied receipts report all 508 exact checks and all 94 diagnostics
as passing. Near a subtraction point the numerical script uses a finite
Taylor series to avoid catastrophic cancellation. That implementation
detail reinforces why the diagnostics are not a replacement for the
analytic Taylor-remainder arguments in the article. A displayed zero
floating-point residual means agreement at the working precision, not
an independently certified exact equality.

\subsection{Source and artifact provenance}

The source manifest records the observed reference commit, the targeted
chapter blob, the scope of the other reads, and the names and Git blob
identifiers of the six incoming archives that were listed but not
opened. It does not assert that all moving-branch reads were a single
atomic snapshot. The targeted mathematical source was re-read at the
reference commit, eliminating that ambiguity for the theorem being
extended.

The package includes the complete \LaTeX{} source, its compiled PDF, an
integration fragment extracted from the mathematical sections, the two
verification scripts, their JSON receipts, a status ledger, build
instructions, and SHA-256 checksums. Build by-products are excluded from
the distribution. The PDF was rendered and visually inspected as part of
package preparation; the separate artifact inspection record documents
that review. No proof-assistant or external peer-review status is inferred
from the artifact checks.

\begin{thebibliography}{99}
\raggedright
\bibitem{repo-mixed}
ProveIt Contributors,
\emph{Unequal twists through beta mixtures and finite resolvent identities},
\texttt{07-mixed-twist-identities.tex}, in the manuscript chapters directory,
reference commit \texttt{446dca9bd65c20e409d1d2561d4ffad408103e2a}.
\href{https://github.com/VladimirReshetnikov/ProveIt/blob/446dca9bd65c20e409d1d2561d4ffad408103e2a/Analysis/Polylogarithms/docs/manuscript/chapters/07-mixed-twist-identities.tex}{Pinned GitHub source}.

\bibitem{repo-readme}
ProveIt Contributors,
\emph{Polylogarithms and their Arithmetic Bridges}, manuscript README,
inspected during this continuation. The moving-branch reading is
recorded in the provenance manifest.
\href{https://github.com/VladimirReshetnikov/ProveIt/tree/main/Analysis/Polylogarithms/docs/manuscript}{GitHub manuscript directory and README}.

\bibitem{repo-harmonic}
ProveIt Contributors,
\emph{Gamma coefficients prove every harmonic-order identity},
\texttt{07-harmonic-gamma-coefficients.tex}, inspected source blob
\texttt{1b18b485517cd551f6a938b1b9736cea942eacf5}.
\href{https://github.com/VladimirReshetnikov/ProveIt/blob/main/Analysis/Polylogarithms/docs/manuscript/chapters/07-harmonic-gamma-coefficients.tex}{GitHub harmonic-coefficient section}.

\bibitem{dlmf-beta}
NIST Digital Library of Mathematical Functions,
\S5.12, \emph{Beta Function}.
\url{https://dlmf.nist.gov/5.12}.

\bibitem{dlmf-gamma}
NIST Digital Library of Mathematical Functions,
\S5.5, \emph{Functional Relations} for the Gamma function.
\url{https://dlmf.nist.gov/5.5}.

\bibitem{dlmf-polygamma}
NIST Digital Library of Mathematical Functions,
\S5.15, \emph{Polygamma Functions}.
\url{https://dlmf.nist.gov/5.15}.

\bibitem{dlmf-polylog}
NIST Digital Library of Mathematical Functions,
\S25.12, \emph{Polylogarithms}.
\url{https://dlmf.nist.gov/25.12}.

\bibitem{dlmf-hurwitz}
NIST Digital Library of Mathematical Functions,
\S25.11, \emph{Hurwitz Zeta Function}.
\url{https://dlmf.nist.gov/25.11}.

\bibitem{weldon}
H.~A. Weldon,
\emph{Mishaps with Feynman parametrization at finite temperature},
Physical Review D \textbf{47} (1993), 594--600.
\href{https://doi.org/10.1103/PhysRevD.47.594}{doi:10.1103/PhysRevD.47.594}.
Abstract consulted for context; no theorem in this article depends on
an uninspected derivation in that paper.

\bibitem{bedaque-das}
P.~F. Bedaque and A.~Das,
\emph{Feynman parametrization and the degenerate electron gas},
Physical Review D \textbf{47} (1993), 601--607;
\href{https://arxiv.org/abs/hep-ph/9209273}{arXiv:hep-ph/9209273}.
\href{https://doi.org/10.1103/PhysRevD.47.601}{doi:10.1103/PhysRevD.47.601}.
Abstract consulted for context.
\end{thebibliography}
\end{document}
